We distinguish ten scenarios by what the observer tries to infer and which
part of the exchange (the transcript) is available.

\smallskip\noindent\textbf{Location and route.}
One packet gives an estimate of the current GPS position (S1). A longer
sequence may reveal a stopping place (S2) or the route travelled, position by
position (S3). For S2 the available proxy is the mean position in the final
minute. The full session also lets us test inference of the starting point
(S9) and endpoint (S10). These are our five main localization tasks.

\smallskip\noindent\textbf{Linkage and prediction.}
Could two sessions belong to the same person or vehicle (S4)? The data
identify route families, so the test asks whether two sessions share a
family. For prediction we distinguish two decisions on the road network.
S5 asks which of two outgoing edges the user will take. S6 asks which of
two destinations the user is heading for once a turn has been observed.

\smallskip\noindent\textbf{Query content and companions.}
For S7 we hold the trip fixed and change the request on the device to see
whether its content changes the transcript. S8 considers another source of
information available to the provider. Given two time-aligned query streams,
how much does observing a companion help locate the user?
